% 第6章 6.6.3：单位权三点路径上多项式滤波的依赖范围。H=I-L/3。
% f=e_1；Hf=(2/3,1/3,0)；H^2 f=(5/9,1/3,1/9)；(L^2-3L)e_1=(-1,0,1)。
% 蓝色：此前已非零；绿色：本轮新到达且非零；虚线白底：在两跳范围内但系数相消为零。
\begin{tikzpicture}[x=1mm,y=1mm,font=\figurefont\footnotesize,text=FigureInk,
  vertex/.style={circle,draw=FigureStroke,fill=white,line width=.8pt,minimum size=6.5mm,inner sep=0pt},
  old/.style={vertex,fill=FigureBlueFill},
  new/.style={vertex,fill=FigureGreenFill,draw=FigureGreen},
  cancel/.style={vertex,dashed,draw=FigureRed},
  edge/.style={draw=FigureStroke,line width=.8pt},
  val/.style={text=FigureBlue},
  note/.style={font=\FigureNoteFont\fontsize{7.5}{9}\selectfont,text=FigureMuted}]
\foreach \y/\op/\sa/\sb/\sc/\va/\vb/\vc/\rem in {
  39/{\bm f=\bm e_1}/old/vertex/vertex/1/0/0/{输入},
  27/{\bm H\bm f}/old/new/vertex/{2/3}/{1/3}/0/{一次乘法：至多一跳},
  15/{\bm H^2\bm f}/old/old/new/{5/9}/{1/3}/{1/9}/{两次乘法：至多两跳},
  3/{(\bm L^2-3\bm L)\bm e_1}/old/cancel/new/{-1}/0/1/{顶点2在范围内，但贡献相消}}{
  \node[anchor=east] at (22,\y) {$\op$};
  \node[\sa] (a) at (32,\y) {$1$};
  \node[\sb] (b) at (50,\y) {$2$};
  \node[\sc] (c) at (68,\y) {$3$};
  \draw[edge] (a)--(b)--(c);
  \node[val,anchor=west] at (35.5,\y+4.3) {$\va$};
  \node[val,anchor=west] at (53.5,\y+4.3) {$\vb$};
  \node[val,anchor=west] at (71.5,\y+4.3) {$\vc$};
  \node[note,anchor=west] at (78,\y) {\rem};
}
\end{tikzpicture}
